\documentclass[12pt,a4paper]{article}
\usepackage[utf8]{inputenc}
\usepackage[portuguese]{babel}
\usepackage{geometry}
\geometry{margin=1in}
\usepackage{graphicx}
\usepackage{listings}
\usepackage{xcolor}
\usepackage{hyperref}
\usepackage{tikz}
\usetikzlibrary{shapes,arrows,positioning}

\lstset{
    language=Python,
    basicstyle=\ttfamily\small,
    breaklines=true,
    keywordstyle=\color{blue},
    commentstyle=\color{gray},
    stringstyle=\color{red},
    showstringspaces=false,
    backgroundcolor=\color{lightgray!20},
    frame=single
}

\title{\textbf{Documentação Completa do Sistema RAG} \\
        \large Explicação Detalhada do Código e do Pipeline}
\author{Documentação do Projeto Prático}
\date{Maio de 2026}

\begin{document}

\maketitle

\tableofcontents
\newpage

% ==================== SEÇÃO 1: INTRODUÇÃO ====================
\section{Introdução ao RAG}

\subsection{O que é RAG?}

RAG significa \textbf{Retrieval-Augmented Generation} (Geração Aumentada por Recuperação).

É uma arquitetura de sistema que combina duas coisas principais:

\begin{enumerate}
    \item \textbf{Recuperação}: Buscar informações relevantes em um banco de dados
    \item \textbf{Geração}: Usar um modelo de linguagem para gerar respostas baseadas nessas informações
\end{enumerate}

\subsection{Por que usar RAG?}

Modelos de linguagem como ChatGPT têm conhecimento geral, mas:

\begin{itemize}
    \item Não sabem sobre documentos específicos do usuário
    \item Podem gerar informações incorretas (``alucinações'')
    \item Não têm acesso a dados privados
\end{itemize}

RAG resolve isso:
\begin{enumerate}
    \item O usuário envia seus documentos (PDFs, TXTs)
    \item O sistema indexa esses documentos
    \item Quando o usuário pergunta algo, o sistema busca as partes relevantes dos documentos
    \item O modelo gera uma resposta baseada APENAS nessas partes recuperadas
\end{enumerate}

\subsection{Vantagens}

\begin{itemize}
    \item \textbf{Precisão}: Respostas baseadas em informações reais do usuário
    \item \textbf{Confiabilidade}: Menos alucinações porque há contexto real
    \item \textbf{Privacidade}: Dados do usuário não saem do seu computador
    \item \textbf{Atualização fácil}: Adicione novos documentos quando necessário
\end{itemize}

% ==================== SEÇÃO 2: PIPELINE RAG ====================
\section{O Pipeline RAG - Fluxo Completo}

\subsection{Visão Geral do Pipeline}

O sistema RAG funciona em dois momentos:

\begin{enumerate}
    \item \textbf{Fase 1 - Indexação (Preparação)}: Processa e armazena documentos
    \item \textbf{Fase 2 - Consulta (Uso)}: Responde perguntas usando os documentos armazenados
\end{enumerate}

\subsection{Diagrama do Pipeline}

\begin{tikzpicture}[node distance=1.5cm]
    % Fase de Indexação
    \node (doc) [rectangle, draw, minimum width=2cm, minimum height=1cm] {Documentos\\(PDF, TXT)};
    \node (load) [rectangle, draw, right of=doc] {Carregar\\Texto};
    \node (chunk) [rectangle, draw, right of=load] {Dividir em\\Segmentos};
    \node (emb) [rectangle, draw, right of=chunk] {Gerar\\Embeddings};
    \node (store) [rectangle, draw, right of=emb] {Armazenar\\no Índice};
    
    \draw[->] (doc) -- (load);
    \draw[->] (load) -- (chunk);
    \draw[->] (chunk) -- (emb);
    \draw[->] (emb) -- (store);
    
    % Fase de Consulta
    \node (q) [rectangle, draw, below of=doc, yshift=-2cm] {Pergunta\\do Usuário};
    \node (qemb) [rectangle, draw, right of=q] {Gerar Embedding\\da Pergunta};
    \node (search) [rectangle, draw, right of=qemb] {Buscar\\Segmentos Similares};
    \node (build) [rectangle, draw, right of=search] {Construir\\Prompt};
    \node (llm) [rectangle, draw, right of=build] {LLM Gera\\Resposta};
    \node (show) [rectangle, draw, right of=llm] {Mostrar\\Resposta};
    
    \draw[->] (q) -- (qemb);
    \draw[->] (qemb) -- (search);
    \draw[->] (search) -- (build);
    \draw[->] (build) -- (llm);
    \draw[->] (llm) -- (show);
    
    % Conexão entre fase 1 e 2
    \draw[dashed] (store) -- (search);
    
    \node at (5, 0.5) {\textbf{FASE 1: Indexação}};
    \node at (5, -2.5) {\textbf{FASE 2: Consulta}};
\end{tikzpicture}

\subsection{Detalhes de Cada Etapa}

\subsubsection{Etapa 1: Ingestão de Dados}

O sistema carrega documentos em formatos PDF ou TXT.

\textbf{Por que múltiplos formatos?}
\begin{itemize}
    \item PDF: Comum para documentos formais, livros
    \item TXT: Simples e universal
\end{itemize}

\subsubsection{Etapa 2: Segmentação (Chunking)}

Documentos são divididos em pequenos pedaços chamados \textbf{chunks}.

\textbf{Por que dividir?}
\begin{itemize}
    \item Modelos têm limite de tokens (palavras)
    \item Chunks pequenos = busca mais precisa
    \item Chunks com overlap = melhor contexto
\end{itemize}

Exemplo: Um documento com 10.000 palavras é dividido em 25 chunks de 400 palavras cada, com 80 palavras de sobreposição.

\subsubsection{Etapa 3: Geração de Embeddings}

Cada chunk é convertido em um \textbf{vetor numérico} (embedding).

Um embedding é como uma ``impressão digital'' do texto que captura seu significado.

\textbf{Exemplo simplificado:}
\begin{itemize}
    \item Texto 1: ``O gato subiu no telhado'' → Vetor [0.2, 0.8, -0.1, 0.5, ...]
    \item Texto 2: ``O felino está no teto'' → Vetor [0.22, 0.79, -0.09, 0.48, ...]
\end{itemize}

Vetores similares = textos com significados similares.

\subsubsection{Etapa 4: Armazenamento em Banco Vetorial}

Os embeddings são armazenados em um \textbf{banco de dados vetorial} (FAISS).

FAISS é rápido em buscas por similaridade: ``Qual embedding é mais próximo deste?''

\subsubsection{Etapa 5: Recuperação}

Quando o usuário faz uma pergunta:
\begin{enumerate}
    \item Gera-se o embedding da pergunta
    \item Busca-se os K embeddings mais similares no banco
    \item Retornam-se os chunks correspondentes (top-k)
\end{enumerate}

\subsubsection{Etapa 6: Construção do Prompt}

Os chunks recuperados são inclusos em um prompt estruturado:

\begin{lstlisting}
Você é um sistema RAG.
Responda SOMENTE usando o contexto fornecido.
NÃO invente informações.

Contexto:
[Aqui vão os chunks recuperados]

Pergunta:
[Aqui vai a pergunta do usuário]

Responda de forma objetiva e curta.
\end{lstlisting}

\subsubsection{Etapa 7: Geração de Resposta}

Um LLM (Large Language Model) como Phi-3 processa o prompt e gera a resposta.

A resposta é baseada APENAS nos chunks recuperados, não em conhecimento geral.

% ==================== SEÇÃO 3: ESTRUTURA DO CÓDIGO ====================
\section{Estrutura do Código}

\subsection{Organização dos Arquivos}

\begin{lstlisting}
ai_/
├── app.py                    # Arquivo principal
├── src/
│   ├── loaders.py            # Carrega PDFs e TXTs
│   ├── chunking.py           # Divide texto em chunks
│   ├── embeddings.py         # Gera embeddings
│   ├── vector_store.py       # Gerencia banco vetorial
│   ├── llm_api.py            # Integração com LLM
│   └── rag_pipeline.py       # Orquestra o pipeline
├── data/                      # Armazena documentos e índices
└── experiment_plan.md        # Plano de testes
\end{lstlisting}

\subsection{Como os Arquivos se Conectam}

\begin{tikzpicture}[node distance=2cm]
    \node (app) [rectangle, draw] {app.py};
    \node (load) [rectangle, draw, below left of=app] {loaders.py};
    \node (chunk) [rectangle, draw, below of=app] {chunking.py};
    \node (emb) [rectangle, draw, below right of=app] {embeddings.py};
    \node (vs) [rectangle, draw, below of=emb] {vector\_store.py};
    \node (llm) [rectangle, draw, below of=chunk] {llm\_api.py};
    \node (rag) [rectangle, draw, below of=vs] {rag\_pipeline.py};
    
    \draw[->] (app) -- (load);
    \draw[->] (app) -- (chunk);
    \draw[->] (app) -- (vs);
    \draw[->] (chunk) -- (emb);
    \draw[->] (emb) -- (vs);
    \draw[->] (app) -- (rag);
    \draw[->] (rag) -- (vs);
    \draw[->] (rag) -- (llm);
\end{tikzpicture}

% ==================== SEÇÃO 4: CADA ARQUIVO ====================
\section{Explicação Detalhada de Cada Arquivo}

\subsection{app.py - O Maestro}

Este é o arquivo principal que controla tudo.

\subsubsection{O que ele faz?}

\begin{enumerate}
    \item Carrega documentos da pasta \texttt{data/}
    \item Processa e indexa os documentos (primeira vez)
    \item Carrega índice existente (próximas vezes)
    \item Inicia um loop interativo para perguntas
    \item Mantém histórico de conversa
\end{enumerate}

\subsubsection{Variáveis de Configuração}

\begin{lstlisting}
CHUNK_SIZE = 400    # Tamanho de cada chunk em palavras
OVERLAP = 80        # Sobreposição entre chunks
TOP_K = 4           # Quantos chunks recuperar
\end{lstlisting}

\textbf{O que significa cada uma?}

\begin{itemize}
    \item \texttt{CHUNK\_SIZE = 400}: Cada pedaço do texto tem 400 palavras
    \item \texttt{OVERLAP = 80}: Os pedaços se sobrepõem em 80 palavras
    \item \texttt{TOP\_K = 4}: Quando o usuário pergunta, busca os 4 chunks mais relevantes
\end{itemize}

\subsubsection{Fluxo Principal}

\begin{enumerate}
    \item \textbf{Carregamento}: Tenta carregar índice antigo
    \item \textbf{Se não existir}: Carrega documentos, divide em chunks, gera embeddings
    \item \textbf{Se existir}: Pula direto para perguntas
    \item \textbf{Loop de Perguntas}: Usuário digita pergunta → sistema responde
    \item \textbf{Histórico}: Cada pergunta e resposta é guardada
\end{enumerate}

\subsubsection{Código Simplificado}

\begin{lstlisting}
# 1. Tentar carregar índice antigo
if load_persisted_store():
    print("Índice carregado!")
else:
    # 2. Se não existir, indexar documentos
    for arquivo em data/:
        texto = carregar_documento(arquivo)
        chunks = dividir_em_chunks(texto)
        gerar_embeddings(chunks)
        armazenar(chunks)

# 3. Loop de perguntas
while True:
    pergunta = input("Digite sua pergunta: ")
    resposta, chunks = answer_question(pergunta, top_k=4)
    print("Resposta:", resposta)
    historico.append({pergunta, resposta})
\end{lstlisting}

% ---------- loaders.py ----------
\subsection{loaders.py - O Leitor}

Este arquivo extrai texto de documentos.

\subsubsection{O que ele faz?}

\begin{enumerate}
    \item Carrega PDFs usando a biblioteca \texttt{fitz}
    \item Carrega TXTs usando Python nativo
    \item Limpa o texto (remove lixo, normaliza espaços)
    \item Retorna texto pronto para processar
\end{enumerate}

\subsubsection{Funções Principais}

\paragraph{load\_txt(file\_path)}
\begin{itemize}
    \item Entrada: caminho do arquivo TXT
    \item Saída: conteúdo do texto
    \item Exemplo: \texttt{``Olá mundo''} (simples)
\end{itemize}

\paragraph{load\_pdf(file\_path)}
\begin{itemize}
    \item Entrada: caminho do arquivo PDF
    \item Processo: Abre o PDF, extrai texto de cada página
    \item Saída: conteúdo de todas as páginas concatenadas
    \item Exemplo: Um PDF com 10 páginas retorna todo o texto junto
\end{itemize}

\paragraph{clean\_text(text)}
\begin{itemize}
    \item Entrada: texto sujo
    \item Limpeza:
    \begin{itemize}
        \item Remove caracteres nulos (\texttt{\textbackslash x00})
        \item Remove marcas de ordem de byte (\texttt{\textbackslash ufeff})
        \item Normaliza quebras de linha
        \item Remove espaços extras
    \end{itemize}
    \item Saída: texto limpo
\end{itemize}

\paragraph{load\_document(file\_path)}
\begin{itemize}
    \item Função principal que coordena as outras
    \item Verifica o tipo de arquivo (.pdf ou .txt)
    \item Chama a função apropriada
    \item Retorna texto limpo
\end{itemize}

\subsubsection{Exemplo de Uso}

\begin{lstlisting}
from src.loaders import load_document

# Carregar um PDF
texto = load_document("documento.pdf")
print(texto)  # Retorna o texto do PDF

# Carregar um TXT
texto = load_document("arquivo.txt")
print(texto)  # Retorna o conteúdo do TXT
\end{lstlisting}

% ---------- chunking.py ----------
\subsection{chunking.py - O Divisor}

Este arquivo divide texto em pedaços menores.

\subsubsection{Por que dividir?}

\begin{itemize}
    \item Modelos têm limite de tokens (palavras)
    \item Phi-3 pode processar ~4000 tokens
    \item Um documento pode ter 100.000+ palavras
    \item Solução: Dividir em chunks de 400 palavras
\end{itemize}

\subsubsection{Como funciona?}

\begin{enumerate}
    \item Divide o texto em palavras
    \item Agrupa 400 palavras em um chunk
    \item Quando atinge 400, salva o chunk
    \item Pega as últimas 80 palavras e começa o próximo (overlap)
    \item Continua até o fim do texto
\end{enumerate}

\subsubsection{Visualização}

\begin{verbatim}
Texto: [palavra1] [palavra2] ... [palavra400] [palavra401] ... [palavra480]
                                    ↓
Chunk 1: [palavra1 até palavra400]
Overlap:                              [palavra321 até palavra400]
                                                  ↓
Chunk 2: [palavra321 até palavra720]
Overlap:                                          [palavra641 até palavra720]
                                                              ↓
Chunk 3: [palavra641 até palavra1040]
... e assim por diante
\end{verbatim}

\subsubsection{Função: chunk\_text}

\begin{lstlisting}
def chunk_text(text, chunk_size=400, overlap=80):
    """
    Divide um texto em chunks com sobreposição.
    
    Parâmetros:
    - text: o texto a dividir
    - chunk_size: tamanho de cada chunk em palavras
    - overlap: quantas palavras da sobreposição
    
    Retorna:
    - Lista de chunks (strings)
    """
    
    # 1. Divide em palavras
    palavras = text.split()
    
    # 2. Cria lista vazia para chunks
    chunks = []
    chunk_atual = []
    
    # 3. Loop por cada palavra
    for palavra in palavras:
        chunk_atual.append(palavra)
        
        # 4. Se atingiu o tamanho
        if len(chunk_atual) >= chunk_size:
            # 5. Salva o chunk
            chunks.append(" ".join(chunk_atual))
            
            # 6. Pega as últimas N palavras para overlap
            chunk_atual = chunk_atual[-overlap:]
    
    # 7. Se sobrou algo, salva
    if chunk_atual:
        chunks.append(" ".join(chunk_atual))
    
    return chunks
\end{lstlisting}

\subsubsection{Exemplo}

\begin{lstlisting}
texto = "palavra1 palavra2 palavra3 ... palavra500"
chunks = chunk_text(texto, chunk_size=10, overlap=2)

# Resultado:
# chunks[0] = "palavra1 palavra2 ... palavra10"
# chunks[1] = "palavra9 palavra10 palavra11 ... palavra20"
# chunks[2] = "palavra19 palavra20 ... palavra30"
# ... e assim por diante
\end{lstlisting}

% ---------- embeddings.py ----------
\subsection{embeddings.py - O Vetorizador}

Este arquivo transforma texto em vetores numéricos.

\subsubsection{O que é um Embedding?}

Um embedding é uma representação numérica do significado de um texto.

\textbf{Analogia:}
\begin{itemize}
    \item Texto é como uma descrição verbal de uma pessoa
    \item Embedding é como uma foto numérica dessa pessoa
    \item Pessoas similares = fotos similares
    \item Textos similares = embeddings similares
\end{itemize}

\subsubsection{Como funciona?}

\begin{enumerate}
    \item Usamos um modelo pré-treinado (all-MiniLM-L6-v2)
    \item Passamos o texto para o modelo
    \item O modelo retorna um vetor de 384 números
    \item Cada número representa uma ``dimensão'' do significado
\end{enumerate}

\subsubsection{Função: generate\_embedding}

\begin{lstlisting}
import torch
from sentence_transformers import SentenceTransformer

# Detecta se GPU está disponível
device = "cuda" if torch.cuda.is_available() else "cpu"

# Carrega o modelo
model = SentenceTransformer("all-MiniLM-L6-v2", device=device)

def generate_embedding(text):
    """
    Gera embedding para um texto.
    
    Entrada: um texto (string)
    Saída: um vetor de 384 números
    """
    embedding = model.encode(text, convert_to_tensor=False)
    return embedding.tolist()
\end{lstlisting}

\subsubsection{Exemplo Prático}

\begin{lstlisting}
# Dois textos similares
texto1 = "O gato subiu no telhado"
texto2 = "O felino está no teto"

emb1 = generate_embedding(texto1)  # [0.2, 0.8, -0.1, 0.5, ...]
emb2 = generate_embedding(texto2)  # [0.22, 0.79, -0.09, 0.48, ...]

# Os embeddings são muito parecidos! (números próximos)

# Texto diferente
texto3 = "Amanhã é segunda-feira"
emb3 = generate_embedding(texto3)  # [-0.5, -0.2, 0.9, ...]

# Este embedding é bem diferente dos dois primeiros
\end{lstlisting}

\subsubsection{GPU vs CPU}

\begin{itemize}
    \item \textbf{GPU}: Muito mais rápido (segundos)
    \item \textbf{CPU}: Lento (minutos)
\end{itemize}

O código detecta automaticamente se há GPU disponível:
\begin{itemize}
    \item Se sim: Usa GPU (muito mais rápido)
    \item Se não: Usa CPU (mais lento)
\end{itemize}

% ---------- vector_store.py ----------
\subsection{vector\_store.py - O Armazenador}

Este arquivo gerencia o banco de dados vetorial.

\subsubsection{O que é um Banco Vetorial?}

Um banco de dados otimizado para armazenar e buscar embeddings rapidamente.

\textbf{Analogia:}
\begin{itemize}
    \item Banco normal: ``Encontre o cliente com nome João''
    \item Banco vetorial: ``Encontre os 4 embeddings mais próximos deste''
\end{itemize}

\subsubsection{Por que FAISS?}

\begin{itemize}
    \item Muito rápido mesmo com milhões de vetores
    \item Usa algoritmos sofisticados para busca eficiente
    \item Desenvolvido pelo Facebook (Meta)
\end{itemize}

\subsubsection{Função: add\_chunks}

\begin{lstlisting}
def add_chunks(chunks):
    """
    Adiciona chunks ao banco vetorial.
    
    Processo:
    1. Para cada chunk
    2. Gera o embedding
    3. Armazena no FAISS
    4. Guarda o texto em documents
    5. Salva em arquivo
    """
    
    for chunk in chunks:
        # 1. Gera embedding
        embedding = generate_embedding(chunk)
        
        # 2. Converte para formato do FAISS
        vector = np.array([embedding]).astype("float32")
        
        # 3. Adiciona ao índice
        index.add(vector)
        
        # 4. Guarda o texto
        documents.append(chunk)
    
    # 5. Salva em arquivo (persistência)
    save_persisted_store()
\end{lstlisting}

\subsubsection{Função: search\_chunks}

\begin{lstlisting}
def search_chunks(query, top_k=4):
    """
    Busca os K chunks mais similares a uma query.
    
    Processo:
    1. Gera embedding da pergunta
    2. Busca no FAISS pelos K mais próximos
    3. Retorna os textos desses chunks
    """
    
    # 1. Gera embedding da pergunta
    query_embedding = generate_embedding(query)
    query_vector = np.array([query_embedding]).astype("float32")
    
    # 2. Busca os K mais próximos
    distances, indices = index.search(query_vector, top_k)
    
    # 3. Pega os textos
    results = []
    for idx in indices[0]:
        results.append(documents[idx])
    
    return results
\end{lstlisting}

\subsubsection{Persistência}

\begin{lstlisting}
def load_persisted_store():
    """Carrega índice salvo em arquivo"""
    if not vector_index.faiss.exists():
        return False  # Não existe ainda
    
    # Carrega do arquivo
    documents = json.load("documents.json")
    index = faiss.read_index("vector_index.faiss")
    
    return True  # Sucesso

def save_persisted_store():
    """Salva índice em arquivo"""
    # Salva os documents em JSON
    json.dump(documents, "documents.json")
    
    # Salva o índice FAISS
    faiss.write_index(index, "vector_index.faiss")
\end{lstlisting}

\textbf{Por que persistência?}
\begin{itemize}
    \item Não precisa reindexar toda vez
    \item Muito mais rápido na segunda execução
    \item Índice é reutilizável
\end{itemize}

% ---------- llm_api.py ----------
\subsection{llm\_api.py - O Gerador}

Este arquivo faz o LLM gerar respostas.

\subsubsection{O que é um LLM?}

\begin{itemize}
    \item LLM = Large Language Model
    \item É um modelo de IA treinado com bilhões de palavras
    \item Consegue entender e gerar texto natural
    \item Exemplos: GPT, ChatGPT, Phi-3
\end{itemize}

\subsubsection{Variáveis de Configuração}

\begin{lstlisting}
MODEL_ID = "microsoft/Phi-3-mini-4k-instruct"
TEMPERATURE = 0.3
\end{lstlisting}

\begin{itemize}
    \item \texttt{MODEL\_ID}: Qual modelo usar
    \item \texttt{TEMPERATURE}: Criatividade vs. Precisão
\end{itemize}

\subsubsection{O que é Temperature?}

\begin{itemize}
    \item \textbf{Temperatura = 0.0}: Sempre escolhe a palavra mais provável (preciso)
    \item \textbf{Temperatura = 0.3}: Levemente criativo (bom para RAG)
    \item \textbf{Temperatura = 0.7}: Muito criativo (risco de alucinação)
    \item \textbf{Temperatura = 1.0}: Muito aleatório (impreciso)
\end{itemize}

Para RAG, usar temperatura baixa (0.0 a 0.3) é melhor.

\subsubsection{Função: ask\_llm}

\begin{lstlisting}
def ask_llm(context, question, history=None):
    """
    Gera resposta do LLM.
    
    Parâmetros:
    - context: chunks recuperados
    - question: pergunta do usuário
    - history: perguntas/respostas anteriores (para chat contextual)
    
    Processo:
    1. Monta o sistema prompt (instruções)
    2. Combina context + question
    3. Passa para o LLM
    4. Retorna a resposta
    """
    
    # 1. Sistema de instruções
    messages = [
        {
            "role": "system",
            "content": (
                "Você é um sistema RAG.\n"
                "Responda SOMENTE usando o contexto fornecido.\n"
                "NÃO invente informações.\n"
                "..."
            )
        },
        {
            "role": "user",
            "content": (
                f"Contexto:\n{context}\n\n"
                f"Pergunta:\n{question}"
            )
        }
    ]
    
    # 2. Converte para formato do modelo
    prompt = tokenizer.apply_chat_template(
        messages,
        tokenize=False,
        add_generation_prompt=True
    )
    
    # 3. Tokeniza (converte em números)
    inputs = tokenizer(prompt, return_tensors="pt")
    
    # 4. Move para GPU se disponível
    inputs = inputs.to(device)
    
    # 5. Gera resposta
    output_ids = model.generate(
        **inputs,
        max_new_tokens=120,
        temperature=TEMPERATURE,
        do_sample=True,
        top_p=0.85,
        repetition_penalty=1.15
    )
    
    # 6. Decodifica (converte de volta para texto)
    response = tokenizer.decode(
        output_ids[0],
        skip_special_tokens=True
    )
    
    return response
\end{lstlisting}

\subsubsection{GPU vs CPU}

Código detecta automaticamente:

\begin{lstlisting}
if torch.cuda.is_available():
    model_kwargs["torch_dtype"] = torch.float16  # GPU: mais rápido
    print("[INFO] GPU detectada!")
else:
    model_kwargs["torch_dtype"] = torch.float32  # CPU: mais lento
    print("[WARNING] GPU não detectada!")
\end{lstlisting}

% ---------- rag_pipeline.py ----------
\subsection{rag\_pipeline.py - O Orquestrador}

Este arquivo coordena todo o pipeline.

\subsubsection{O que faz?}

Conecta três partes:
\begin{enumerate}
    \item Recuperação (search\_chunks)
    \item Prompt (combinação de contexto)
    \item Geração (ask\_llm)
\end{enumerate}

\subsubsection{Função: answer\_question}

\begin{lstlisting}
def answer_question(question, top_k=4, history=None):
    """
    Responde uma pergunta usando o pipeline RAG.
    
    Processo:
    1. Recupera chunks similares
    2. Combina em contexto
    3. Passa para LLM
    4. Retorna resposta + chunks (para transparência)
    """
    
    # 1. Recupera
    retrieved_chunks = search_chunks(question, top_k)
    
    # 2. Combina
    context = "\n\n".join(retrieved_chunks)
    
    # 3. Gera
    response = ask_llm(
        context,
        question,
        history=history
    )
    
    # 4. Retorna
    return response, retrieved_chunks
\end{lstlisting}

% ==================== SEÇÃO 5: FLUXO COMPLETO ====================
\section{Fluxo Completo do Sistema}

\subsection{Primeira Execução}

\begin{enumerate}
    \item \textbf{app.py inicia}
    \item \textbf{Tenta carregar índice} → Não existe
    \item \textbf{Carrega documentos}
    \begin{itemize}
        \item \texttt{loaders.py} carrega PDFs/TXTs
        \item \texttt{chunking.py} divide em chunks
    \end{itemize}
    \item \textbf{Gera embeddings}
    \begin{itemize}
        \item \texttt{embeddings.py} vetoriza cada chunk
        \item \texttt{vector\_store.py} armazena no FAISS
    \end{itemize}
    \item \textbf{Salva índice} em arquivo
    \item \textbf{Aguarda perguntas}
\end{enumerate}

\subsection{Próximas Execuções}

\begin{enumerate}
    \item \textbf{app.py inicia}
    \item \textbf{Tenta carregar índice} → Existe!
    \item \textbf{Pula reindexação} (muito mais rápido)
    \item \textbf{Aguarda perguntas}
\end{enumerate}

\subsection{Respondendo uma Pergunta}

\begin{enumerate}
    \item Usuário digita: ``Quais são os personagens?''
    \item \textbf{app.py} chama \texttt{answer\_question()}
    \item \textbf{rag\_pipeline.py}:
    \begin{enumerate}
        \item Chama \texttt{search\_chunks(``Quais são os personagens?'')}
        \item \textbf{vector\_store.py}:
        \begin{enumerate}
            \item \textbf{embeddings.py} gera embedding da pergunta
            \item FAISS busca os 4 chunks mais similares
            \item Retorna os textos
        \end{enumerate}
        \item Combina os chunks em um contexto
        \item Chama \texttt{ask\_llm(context, question)}
        \item \textbf{llm\_api.py}:
        \begin{enumerate}
            \item Monta prompt com instrução + contexto + pergunta
            \item Envia para Phi-3
            \item Phi-3 gera resposta baseada no contexto
            \item Retorna a resposta
        \end{enumerate}
    \end{enumerate}
    \item Mostra resposta ao usuário
    \item Armazena no histórico
\end{enumerate}

% ==================== SEÇÃO 6: PARÂMETROS ====================
\section{Parâmetros e Configurações}

\subsection{Parâmetros Principais}

\subsubsection{CHUNK\_SIZE (tamanho do chunk)}

\begin{itemize}
    \item \textbf{Padrão}: 400
    \item \textbf{Pequeno} (300): Mais chunks, busca mais precisa, mas mais lento
    \item \textbf{Grande} (500): Menos chunks, mais contexto, possível sobrecarregar
\end{itemize}

\subsubsection{OVERLAP (sobreposição)}

\begin{itemize}
    \item \textbf{Padrão}: 80
    \item \textbf{0}: Sem overlap, rápido mas pode perder contexto
    \item \textbf{80}: Com overlap, melhor contexto
    \item \textbf{Até 200}: Overlap maior, muito contexto
\end{itemize}

\subsubsection{TOP\_K (quantos chunks recuperar)}

\begin{itemize}
    \item \textbf{Padrão}: 4
    \item \textbf{2}: Rápido, mas pode faltar contexto
    \item \textbf{4}: Bom balanço
    \item \textbf{6+}: Muito contexto, pode ser redundante
\end{itemize}

\subsubsection{TEMPERATURE (criatividade)}

\begin{itemize}
    \item \textbf{Padrão}: 0.3
    \item \textbf{0.0}: Muito preciso, repetitivo
    \item \textbf{0.3}: Bom para RAG
    \item \textbf{0.7}: Criativo, pode alucinar
\end{itemize}

\subsection{Onde Alterar}

\subsubsection{app.py}

\begin{lstlisting}
# No início do arquivo
CHUNK_SIZE = 400   # Altere aqui
OVERLAP = 80       # Altere aqui
TOP_K = 4          # Altere aqui
\end{lstlisting}

\subsubsection{llm\_api.py}

\begin{lstlisting}
# No início do arquivo
MODEL_ID = "microsoft/Phi-3-mini-4k-instruct"  # Altere aqui
TEMPERATURE = 0.3                                # Altere aqui
\end{lstlisting}

% ==================== SEÇÃO 7: ERROS COMUNS ====================
\section{Erros Comuns e Soluções}

\subsection{``ModuleNotFoundError: No module named 'fitz''}

\textbf{Solução:}
\begin{lstlisting}
pip install PyMuPDF
\end{lstlisting}

\subsection{``CUDA out of memory''}

\textbf{Possíveis causas:}
\begin{itemize}
    \item TOP\_K muito grande
    \item Chunks muito grandes
    \item Modelo carregado errado
\end{itemize}

\textbf{Soluções:}
\begin{enumerate}
    \item Reduzir TOP\_K para 2 ou 3
    \item Reduzir CHUNK\_SIZE para 300
    \item Usar torch.float16 (já automático)
\end{enumerate}

\subsection{Resposta vazia ou ``Erro na geração''}

\textbf{Possíveis causas:}
\begin{itemize}
    \item Nenhum chunk recuperado
    \item LLM travado
    \item Contexto vazio
\end{itemize}

\textbf{Soluções:}
\begin{enumerate}
    \item Verificar se documentos foram carregados
    \item Aumentar TOP\_K
    \item Checar se os documentos contêm informação relevante
\end{enumerate}

\subsection{Sistema muito lento}

\textbf{Causa provável:} Usando CPU em vez de GPU

\textbf{Solução:}
\begin{enumerate}
    \item Instalar CUDA
    \item Verificar: \texttt{python -c "import torch; print(torch.cuda.is\_available())"}
    \item Se retorna \texttt{False}: CUDA não está instalado
\end{enumerate}

% ==================== SEÇÃO 8: RESUMO VISUAL ====================
\section{Resumo Visual do Sistema}

\subsection{Componentes}

\begin{table}[h]
    \centering
    \begin{tabular}{|l|l|l|}
        \hline
        \textbf{Arquivo} & \textbf{Função} & \textbf{Saída} \\
        \hline
        loaders.py & Carrega PDFs/TXTs & Texto limpo \\
        \hline
        chunking.py & Divide em chunks & Lista de chunks \\
        \hline
        embeddings.py & Vetoriza chunks & Vetores (384 números) \\
        \hline
        vector\_store.py & Armazena no FAISS & Índice persistente \\
        \hline
        llm\_api.py & Gera respostas & Texto natural \\
        \hline
        rag\_pipeline.py & Coordena tudo & Resposta final \\
        \hline
        app.py & Interface interativa & Chat com usuário \\
        \hline
    \end{tabular}
\end{table}

\subsection{Fluxo de Dados}

\begin{verbatim}
Documentos (PDF/TXT)
        ↓ [loaders.py]
Texto Limpo
        ↓ [chunking.py]
Chunks de 400 palavras
        ↓ [embeddings.py]
Vetores de 384 dimensões
        ↓ [vector_store.py]
Índice FAISS (arquivo)
        ↓
[Na próxima pergunta]
        ↓ [embeddings.py]
Pergunta vetorizada
        ↓ [vector_store.py]
Top-4 chunks similares
        ↓ [llm_api.py]
Resposta do LLM
        ↓
Usuário vê a resposta
\end{verbatim}

% ==================== SEÇÃO 9: MELHORIAS FUTURAS ====================
\section{Melhorias e Otimizações Futuras}

\subsection{Requisitos Não Implementados (Projeto)}

\subsubsection{1. Interface Gráfica (CRÍTICO)}

\textbf{Atual:} CLI (linha de comando com input/print)

\textbf{Melhorado:} Streamlit ou Dash

\begin{lstlisting}
# Exemplo com Streamlit
import streamlit as st

st.title("Sistema RAG")

uploaded_file = st.file_uploader("Envie um PDF")
if uploaded_file:
    indexar_documento(uploaded_file)

pergunta = st.text_input("Faça uma pergunta")
if pergunta:
    resposta, chunks = answer_question(pergunta)
    st.write("Resposta:", resposta)
    st.write("Chunks recuperados:", chunks)
\end{lstlisting}

\textbf{Benefícios:}
\begin{itemize}
    \item Interface visual atraente
    \item Upload de arquivos fácil
    \item Histórico de conversa visual
    \item Visualização de chunks recuperados
\end{itemize}

\subsubsection{2. APIs Externas (CRÍTICO)}

\textbf{Atual:} Modelos locais (CPU/GPU)
\begin{itemize}
    \item Embeddings: sentence-transformers (local)
    \item LLM: Phi-3 (local)
\end{itemize}

\textbf{Requisito do Projeto:} APIs externas

\textbf{Sugestões:}

\paragraph{Para Embeddings:}
\begin{itemize}
    \item OpenAI Embeddings API
    \item Cohere API
    \item HuggingFace Inference API
\end{itemize}

\paragraph{Para LLM:}
\begin{itemize}
    \item OpenAI GPT-4 / GPT-3.5
    \item Cohere Command
    \item Google PaLM
    \item Anthropic Claude
\end{itemize}

\textbf{Exemplo com OpenAI:}

\begin{lstlisting}
import openai

# Embeddings
response = openai.Embedding.create(
    input="Qual é a cor do céu?",
    model="text-embedding-3-small"
)
embedding = response['data'][0]['embedding']

# LLM
response = openai.ChatCompletion.create(
    model="gpt-3.5-turbo",
    messages=[{"role": "user", "content": prompt}]
)
resposta = response['choices'][0]['message']['content']
\end{lstlisting}

\textbf{Vantagens:}
\begin{itemize}
    \item Não precisa de GPU local
    \item Modelos mais poderosos
    \item Atende requisito do projeto
\end{itemize}

\textbf{Desvantagens:}
\begin{itemize}
    \item Custa dinheiro (API cobrada)
    \item Dados vão para servidor externo
    \item Depende de internet
\end{itemize}

\subsection{Otimizações de Código}

\subsubsection{1. Cache de Embeddings}

\textbf{Problema Atual:} Regera embedding sempre que processa

\textbf{Solução:}

\begin{lstlisting}
from functools import lru_cache

@lru_cache(maxsize=1000)
def generate_embedding(text):
    """Com cache, texto iguais reutilizam embedding"""
    embedding = model.encode(text)
    return embedding.tolist()
\end{lstlisting}

\textbf{Ganho:} 10-100x mais rápido em buscas repetidas

\subsubsection{2. Tratamento de Erros Robusto}

\textbf{Atual:} Mínimo, falhas geralmente travam o programa

\textbf{Melhorado:}

\begin{lstlisting}
def load_document_safe(file_path):
    """Versão segura que trata erros"""
    try:
        if not Path(file_path).exists():
            raise FileNotFoundError(f"Arquivo não encontrado: {file_path}")
        
        text = load_document(file_path)
        
        if not text:
            raise ValueError(f"Arquivo vazio: {file_path}")
        
        return text
    
    except FileNotFoundError as e:
        logger.error(f"Erro de arquivo: {e}")
        return None
    
    except Exception as e:
        logger.error(f"Erro inesperado: {e}")
        return None
\end{lstlisting}

\subsubsection{3. Logging Detalhado}

\textbf{Adicionar:}

\begin{lstlisting}
import logging

logging.basicConfig(
    level=logging.INFO,
    format='%(asctime)s - %(levelname)s - %(message)s'
)
logger = logging.getLogger(__name__)

# Usar em todo o código
logger.info("Iniciando carregamento de documentos")
logger.warning("Índice não encontrado, reindexando")
logger.error("Erro ao gerar embedding")
\end{lstlisting}

\textbf{Benefício:} Rastrear exatamente o que acontece

\subsubsection{4. Validação de Entrada}

\textbf{Adicionar:}

\begin{lstlisting}
def validate_question(question):
    """Valida pergunta antes de processar"""
    
    # Não vazio
    if not question or len(question.strip()) == 0:
        raise ValueError("Pergunta não pode estar vazia")
    
    # Não muito longa
    if len(question) > 5000:
        raise ValueError("Pergunta muito longa (máx 5000 caracteres)")
    
    # Não caracteres inválidos
    if not question.isascii() and not question.isidentifier():
        # Permitir caracteres especiais em português
        pass
    
    return question.strip()

# Usar no app
pergunta_validada = validate_question(pergunta)
\end{lstlisting}

\subsubsection{5. Configuração por Arquivo}

\textbf{Ao invés de:}

\begin{lstlisting}
# No código
CHUNK_SIZE = 400
TOP_K = 4
TEMPERATURE = 0.3
\end{lstlisting}

\textbf{Fazer:}

\begin{lstlisting}
# Arquivo: config.yaml
chunking:
  chunk_size: 400
  overlap: 80

retrieval:
  top_k: 4
  similarity_threshold: 0.5

generation:
  temperature: 0.3
  max_tokens: 120
  model_id: "microsoft/Phi-3-mini-4k-instruct"
\end{lstlisting}

\begin{lstlisting}
# Código: config.py
import yaml

def load_config(config_path="config.yaml"):
    with open(config_path) as f:
        config = yaml.safe_load(f)
    return config

# Usar
config = load_config()
CHUNK_SIZE = config['chunking']['chunk_size']
TOP_K = config['retrieval']['top_k']
\end{lstlisting}

\textbf{Vantagem:} Não precisa editar código para mudar parâmetros

\subsubsection{6. Suporte a Múltiplos Formatos}

\textbf{Atual:} Apenas PDF e TXT

\textbf{Adicionar:}

\begin{itemize}
    \item DOCX (Word)
    \item XLSX (Excel)
    \item PPTX (PowerPoint)
    \item HTML
    \item Markdown
\end{itemize}

\begin{lstlisting}
# Exemplo: DOCX
from docx import Document

def load_docx(file_path):
    doc = Document(file_path)
    text = "\n".join([p.text for p in doc.paragraphs])
    return text

# Adicionar ao loaders.py
if suffix == ".docx":
    text = load_docx(file_path)
\end{lstlisting}

\subsubsection{7. Testes Unitários}

\textbf{Criar:}

\begin{lstlisting}
# tests/test_chunking.py
import pytest
from src.chunking import chunk_text

def test_chunk_size():
    texto = " ".join(["palavra"] * 500)
    chunks = chunk_text(texto, chunk_size=10, overlap=2)
    
    assert len(chunks) > 0
    assert len(chunks[0].split()) >= 10

def test_overlap():
    chunks = chunk_text("a b c d e f g h i j", chunk_size=3, overlap=1)
    # Verificar se há overlap
    assert "d" in chunks[0] or "d" in chunks[1]

# Rodar testes
pytest tests/
\end{lstlisting}

\textbf{Benefício:} Garante que o código funciona

\subsubsection{8. Documentação Inline}

\textbf{Adicionar docstrings em formato Google:}

\begin{lstlisting}
def chunk_text(text, chunk_size=400, overlap=80):
    """Divide texto em chunks com sobreposição.
    
    Args:
        text (str): Texto a dividir
        chunk_size (int): Tamanho de cada chunk em palavras. Default 400.
        overlap (int): Sobreposição em palavras. Default 80.
    
    Returns:
        list: Lista de strings (chunks)
    
    Raises:
        ValueError: Se chunk_size < 0 ou overlap < 0
        TypeError: Se text não é string
    
    Examples:
        >>> chunks = chunk_text("palavra1 palavra2 ... palavra500", 10, 2)
        >>> len(chunks) > 0
        True
    """
    ...
\end{lstlisting}

\subsubsection{9. Monitoramento de Performance}

\textbf{Adicionar métricas:}

\begin{lstlisting}
import time

def measure_time(func):
    def wrapper(*args, **kwargs):
        start = time.time()
        result = func(*args, **kwargs)
        elapsed = time.time() - start
        logger.info(f"{func.__name__} levou {elapsed:.2f}s")
        return result
    return wrapper

@measure_time
def add_chunks(chunks):
    ...

@measure_time
def search_chunks(query, top_k=3):
    ...
\end{lstlisting}

\textbf{Resultado:}
\begin{verbatim}
add_chunks levou 45.32s
search_chunks levou 0.15s
\end{verbatim}

\subsubsection{10. Chat Melhorado}

\textbf{Atual:} Histórico simples

\textbf{Melhorado:}

\begin{itemize}
    \item Salvar conversa em arquivo
    \item Retomar conversa anterior
    \item Feedback do usuário (resposta boa/ruim)
    \item Sugestões automáticas de perguntas
\end{itemize}

\begin{lstlisting}
# Salvar conversa
import json
from datetime import datetime

def save_conversation(history):
    filename = f"conversations/{datetime.now():%Y%m%d_%H%M%S}.json"
    with open(filename, 'w') as f:
        json.dump(history, f, indent=2)

# Carregar conversa anterior
def load_conversation(filename):
    with open(filename) as f:
        return json.load(f)
\end{lstlisting}

\subsubsection{11. Compressão de Contexto}

\textbf{Problema:} Contexto muito longo pode confundir o LLM

\textbf{Solução:}

\begin{lstlisting}
def compress_context(chunks, max_length=2000):
    """Comprime chunks para não exceder tamanho máximo"""
    context = ""
    for chunk in chunks:
        if len(context) + len(chunk) <= max_length:
            context += chunk + "\n\n"
        else:
            break
    return context
\end{lstlisting}

\subsubsection{12. Busca Híbrida}

\textbf{Atual:} Busca apenas por similaridade vetorial

\textbf{Melhorado:} Combinar busca vetorial + busca por palavras-chave

\begin{lstlisting}
def search_hybrid(query, top_k=4):
    """Busca híbrida: vetorial + keyword"""
    
    # 1. Busca vetorial
    vector_results = search_chunks(query, top_k)
    
    # 2. Busca por palavras-chave
    keywords = query.split()
    keyword_results = [
        doc for doc in documents
        if any(kw in doc.lower() for kw in keywords)
    ]
    
    # 3. Combina e ordena
    combined = vector_results + keyword_results
    unique = list(dict.fromkeys(combined))  # Remove duplicatas
    
    return unique[:top_k]
\end{lstlisting}

\subsection{Resumo das Melhorias}

\begin{table}[h]
    \centering
    \begin{tabular}{|l|l|l|}
        \hline
        \textbf{Melhoria} & \textbf{Dificuldade} & \textbf{Impacto} \\
        \hline
        Interface Streamlit & Média & Alto (requisito) \\
        APIs Externas & Alta & Alto (requisito) \\
        Cache de Embeddings & Baixa & Alto (velocidade) \\
        Tratamento de Erros & Média & Médio (robustez) \\
        Logging & Baixa & Médio (debug) \\
        Validação de Entrada & Baixa & Médio (segurança) \\
        Config. por Arquivo & Média & Médio (flexibilidade) \\
        Testes Unitários & Média & Médio (confiabilidade) \\
        Documentação & Baixa & Médio (manutenção) \\
        Chat Melhorado & Média & Baixo (UX) \\
        Busca Híbrida & Alta & Alto (precisão) \\
        Performance & Média & Alto (velocidade) \\
        \hline
    \end{tabular}
\end{table}


\subsection{O que você aprendeu}

\begin{enumerate}
    \item \textbf{RAG é um pipeline}: Recuperação + Geração
    \item \textbf{Cada arquivo tem um propósito}: Divisão de responsabilidades
    \item \textbf{Tudo se conecta}: Dados fluem de um arquivo para outro
    \item \textbf{Parâmetros importam}: Ajuste chunk\_size, top\_k, temperature
    \item \textbf{GPU é muito mais rápido}: Use quando disponível
    \item \textbf{Persistência economiza tempo}: Índice é reutilizável
\end{enumerate}

\subsection{Próximos Passos}

\begin{enumerate}
    \item Execute o sistema com os parâmetros padrão
    \item Teste os experimentos (altere CHUNK\_SIZE, TOP\_K, etc)
    \item Meça o tempo e a qualidade das respostas
    \item Documente os resultados em slides
    \item Implemente a interface Streamlit (requisito do projeto)
\end{enumerate}

\subsection{Dica Final}

\textbf{Comece simples, teste uma coisa por vez.}

Altere um parâmetro, veja o resultado, anote. Depois altere outro. Assim você entende o sistema de verdade.

\end{document}
